\documentclass[11pt,a4paper]{article}

% Packages
\usepackage[utf8]{inputenc}
\usepackage[T1]{fontenc}
\usepackage{amsmath,amssymb,amsthm}
\usepackage{booktabs}
\usepackage{hyperref}
\usepackage[margin=1in]{geometry}
\usepackage{microtype}
\usepackage{natbib}
\usepackage{graphicx}

% Theorem environments
\newtheorem{theorem}{Theorem}
\newtheorem{conjecture}{Conjecture}
\newtheorem{lemma}{Lemma}
\newtheorem{definition}{Definition}
\newtheorem{remark}{Remark}
\newtheorem{observation}{Observation}

\title{Lookahead Depth Tracks Constraint Graph Diameter\\in Random 3-SAT}
\author{Christopher Powers\thanks{Correspondence: \texttt{powers.chr@gmail.com}}}
\date{April 2026}

\begin{document}
\maketitle

\begin{abstract}
We report an empirical connection between $k$-step lookahead depth in DPLL-based SAT solving and the constraint graph diameter of random 3-SAT instances. At each decision point, the solver evaluates every (variable, value) pair by simulating $k$ levels of unit propagation, choosing the assignment with the highest cumulative yield. We measure the largest problem size $n$ at which this achieves a 100\% zero-backtrack rate on hard-core instances (those requiring backtracks under standard heuristics) at clause-to-variable ratio 4.0, and find: $n_{\mathrm{perfect}}(2) = 15$, $n_{\mathrm{perfect}}(3) = 47$, and $n_{\mathrm{perfect}}(4) \geq 100$, with growth factors of $3.13\times$ and ${\geq}2.13\times$. Beam-search lower bounds extend through $k = 7$ ($n \geq 180$). Separately, we measure the constraint graph diameter as $d \approx 0.40 \cdot \log_2 n$, constant across three orders of magnitude ($n = 5$ to $10{,}000$). A linear fit of the required lookahead depth gives $k \approx 1.04 \cdot \log_2 n - 1.94$, closely tracking the diameter. We validate the scaling pattern on SATLIB benchmarks (uf20 through uf100). These results hold at ratio 4.0, above the clustering threshold (${\sim}3.86$) where known rigorous polynomial algorithms fail. We describe a rank-correlation analysis showing that deeper lookahead does not refine but \emph{reshuffles} candidate rankings, and outline open questions about the asymptotic behavior of the zero-backtrack boundary.
\end{abstract}

\noindent\textbf{Keywords:} SAT solving, lookahead heuristics, variable ordering, backtrack-free search, constraint graph diameter, phase transition

\section{Introduction}

\subsection{Background}

The satisfiability problem (SAT) is the canonical NP-complete problem~\citep{cook1971complexity}. DPLL-based solvers~\citep{davis1962machine} search for satisfying assignments by choosing variables, assigning values, propagating unit clauses, and backtracking on contradiction. The choice of variable ordering dramatically affects performance: a good ordering can eliminate backtracking entirely, while a poor one leads to exponential search.

Lookahead solvers~\citep{heule2009lookahead} invest computation at each decision point to choose better variables. Rather than relying on cheap heuristics like Jeroslow--Wang~\citep{jeroslow1990solving}, they simulate multiple levels of propagation to estimate each candidate's downstream consequences. The trade-off is clear: deeper lookahead makes better decisions but costs more per decision.

This raises a natural question: for a given problem size $n$, how deep must the lookahead be to eliminate backtracking entirely? And how does that required depth scale with $n$?

\subsection{Our Contribution}

We present the first systematic measurement of the \emph{zero-backtrack boundary} as a function of lookahead depth $k$ on random 3-SAT. Our main findings:

\begin{enumerate}
    \item \textbf{The required depth tracks constraint graph diameter.} The lookahead depth needed for zero backtracks on all hard-core instances grows as $k \approx 1.04 \cdot \log_2 n$, closely matching the constraint graph diameter $d \approx 0.40 \cdot \log_2 n$. This connection has not, to our knowledge, been previously reported.

    \item \textbf{Three exact data points establish superlinear growth.} The zero-backtrack boundary satisfies $n_{\mathrm{perfect}}(2) = 15$, $n_{\mathrm{perfect}}(3) = 47$, $n_{\mathrm{perfect}}(4) \geq 100$, ruling out a linear relationship $n \propto k$.

    \item \textbf{The pattern holds above the clustering threshold.} All experiments use ratio 4.0, above the barrier at ${\sim}3.86$ where the solution space shatters~\citep{achlioptas2008algorithmic} and where the best known rigorous polynomial algorithm ceases to work~\citep{cojaoghlan2010better}.

    \item \textbf{Deeper lookahead reshuffles, not refines.} Rank-correlation analysis shows that $k{+}1$ does not break ties within $k$'s top candidates but produces a qualitatively different ranking, with correlation dropping precisely at the boundary where $k$ begins to fail.
\end{enumerate}

We emphasize that these are empirical observations on small instances ($n \leq 180$). Whether the scaling law holds asymptotically, and what complexity-theoretic consequences it might have, remain open questions that we discuss but do not resolve.

\section{Methodology}

\subsection{Instance Generation}

We generate random 3-SAT instances using the standard uniform model: for $n$ variables at clause-to-variable ratio $r$, we create $\lfloor n \cdot r \rfloor$ clauses, each containing 3 distinct variables with uniformly random polarity. All experiments use ratio $r = 4.0$, near the satisfiability threshold~\citep{dingsly2015proof}.

Instances are generated from sequential integer seeds using a deterministic xorshift64 PRNG, ensuring full reproducibility. Every reported failure includes its seed for independent verification.

\subsection{Hard Core Definition}

\begin{definition}
An instance is \emph{hard core} if it is satisfiable but requires at least one backtrack under both:
\begin{itemize}
    \item \textbf{Adaptive polarity:} Choose the variable with the most biased polarity, set to the majority direction.
    \item \textbf{Jeroslow--Wang (JW):} Choose the variable maximizing $\sum_c 2^{-|c|}$ over clauses $c$ containing the literal~\citep{jeroslow1990solving}, set to the direction with higher weight.
\end{itemize}
\end{definition}

The hard core fraction grows with $n$ at ratio 4.0: from 4\% at $n = 6$ to 38\% at $n = 30$. We note that these baseline heuristics are not state-of-the-art; a comparison with modern CDCL solvers on the same instances is an important direction for future work (Section~\ref{sec:future}).

\subsection{Solver Implementations}

We developed three C solver implementations:

\subsubsection{Exact Bitwise Solver}

The gold-standard implementation uses 128-bit bitmask clause representation for bitwise unit propagation, popcount-based cardinality checks, and zero heap allocation. The scoring function evaluates every (variable, value) pair at every lookahead level with no pruning.

\subsubsection{Parallel Exact Solver}

An OpenMP-parallelized version distributes top-level scoring across threads, achieving ${\sim}8$--$9\times$ speedup on 16 threads. This extends exact computation to $n = 100$ for $k = 4$.

\subsubsection{Beam-Search Solver}

For larger problems, a beam-search approximation evaluates only the top-$B$ variables (ranked by JW score) at each recursive level, reducing the per-decision cost from $O(n^{2k})$ to $O(B^{k-1} \cdot n)$.

\subsubsection{Calibration}

Beam search systematically underestimates the zero-backtrack boundary. At $k = 3$: exact boundary $n = 47$, beam (width 20) boundary $n = 38$, giving a correction factor of $47/38 = 1.237$. This calibration is measured at a single $k$ value, and all beam-estimated boundaries should be treated as lower bounds.

\subsection{Experimental Protocol}

For each $(k, n)$ pair: generate instances from sequential seeds, filter for hard core, solve with $k$-step lookahead, record zero-backtrack count and wall time. Sample sizes range from 8--50 depending on per-instance compute cost. Time limits: 120 seconds per instance.

\section{Results}

\subsection{Zero-Backtrack Boundaries}

\textbf{$k = 1$:} Single-step lookahead achieves 50--60\% zero-backtrack rate on the hard core with no consistent trend in $n$.

\textbf{$k = 2$:} 100\% zero-backtrack rate through $n = 15$. First failure at $n = 18$ (seed~14, 11~backtracks). Degradation gradual beyond the boundary.

\textbf{$k = 3$:} 100\% through $n = 47$. First failure at $n = 48$ (seed~27, 1506~backtracks). Growth factor: $47/15 = 3.13\times$.

\textbf{$k = 4$:} Parallel exact solver confirms 100\% through $n = 100$ on all hard-core instances. All 14 tested sizes (56--100) achieved zero backtracks. Wall times ranged from 135s ($n = 56$) to 2241s ($n = 95$). Growth factor from $k = 3$: at minimum $100/47 = 2.13\times$.

\textbf{$k = 5$--$7$ (beam estimates):} $k = 5$ beam=6 perfect through $n = 125$; $k = 6$ beam=3 through $n = 160$; $k = 7$ beam=2 through $n = 180$. These are lower bounds; the true boundaries are likely higher, and the decreasing beam widths make them increasingly conservative.

\subsection{The Scaling Law}

\begin{table}[h]
\centering
\begin{tabular}{@{}ccccl@{}}
\toprule
$k$ & $n_{\mathrm{confirmed}}$ & $n/k$ & Growth & Method \\
\midrule
2 & 15 & 7.5  & ---       & Exact \\
3 & 47 & 15.7 & $3.13\times$ & Exact \\
4 & ${\geq}100$ & ${\geq}25$ & ${\geq}2.13\times$ & Exact \\
5 & ${\geq}125$ & ${\geq}25$ & ---       & Beam (lower bound) \\
6 & ${\geq}160$ & ${\geq}27$ & ---       & Beam (lower bound) \\
7 & ${\geq}180$ & ${\geq}26$ & ---       & Beam (lower bound) \\
\bottomrule
\end{tabular}
\caption{Zero-backtrack boundary by lookahead depth. Three exact data points confirm superlinear growth: the $n/k$ ratio is monotonically increasing.}
\label{tab:scaling}
\end{table}

The monotonically increasing $n/k$ ratio (7.5, 15.7, ${\geq}25$) rules out the linear hypothesis $n_{\mathrm{perfect}} = ck$, which would give constant $n/k$. A linear fit across all data points yields:
\begin{equation}\label{eq:fit}
    k \approx 1.04 \cdot \log_2 n - 1.94
\end{equation}

We stress that three exact data points, supplemented by four beam-estimated lower bounds, are a limited basis for fitting a functional form. The growth factor is decelerating ($3.13\times$ to ${\geq}2.13\times$), and we cannot determine from the available data whether it stabilizes above any particular value, continues to decelerate, or eventually collapses.

\subsection{Constraint Graph Diameter}\label{sec:diameter}

The constraint graph connects variables that share at least one clause. We measured its diameter via BFS:

\begin{table}[h]
\centering
\begin{tabular}{@{}cccc@{}}
\toprule
$n$ & Diameter $d$ & $\log_2 n$ & $d / \log_2 n$ \\
\midrule
10     & 2.0 & 3.32  & 0.60 \\
50     & 2.5 & 5.64  & 0.45 \\
100    & 3.0 & 6.64  & 0.45 \\
500    & 3.7 & 8.97  & 0.41 \\
1{,}000  & 4.0 & 9.97  & 0.40 \\
5{,}000  & 4.8 & 12.29 & 0.39 \\
10{,}000 & 5.0 & 13.29 & 0.38 \\
\bottomrule
\end{tabular}
\caption{Constraint graph diameter. The ratio $d/\log_2 n$ stabilizes near 0.40 across three orders of magnitude.}
\label{tab:diameter}
\end{table}

That $d = O(\log n)$ for random 3-SAT at constant ratio follows from the expansion properties of the underlying random hypergraph. The empirical observation is that the required lookahead depth $k$ from~\eqref{eq:fit} closely tracks this diameter: the solver needs to ``see'' approximately 2.5 diameters to eliminate all backtracking.

\begin{observation}
The ratio between the empirically required lookahead depth and the constraint graph diameter is approximately constant: $k / d \approx 2.5$ across all data points where both quantities are measured.
\end{observation}

This suggests a structural explanation: $k$-step lookahead gathers information within graph distance $k$ of the decision variable. When $k$ exceeds the diameter, the scorer sees every constraint in the formula. Between 1 and 2.5 diameters, coverage transitions from partial to complete, and the zero-backtrack rate transitions from partial to perfect.

\subsection{SATLIB Validation}

We validated the scaling pattern on standard SATLIB benchmarks~\citep{satlib}:

\begin{table}[h]
\centering
\begin{tabular}{@{}ccccc@{}}
\toprule
Benchmark & $n$ & Ratio & $k$ for 100\% HC & Previous $k$ \\
\midrule
uf20  & 20  & 4.27 & 3 & $k{=}2$: 98\% \\
uf50  & 50  & 4.36 & 4 & $k{=}3$: 97\% \\
uf100 & 100 & 4.30 & 5 & $k{=}4$: 79\% \\
\bottomrule
\end{tabular}
\caption{SATLIB validation. Each $k$ step catches what the previous misses. SATLIB ratios are slightly above 4.0.}
\label{tab:satlib}
\end{table}

The progression matches our generated instances: the required $k$ for perfect performance on the hard core increases in step with $n$, consistent with the logarithmic fit.

\section{The Reshuffle Mechanism}\label{sec:reshuffle}

To understand \emph{why} additional depth helps, we analyzed how scoring rankings change between $k$ and $k+1$.

\subsection{Rank Correlation}

We computed Spearman rank correlation $\rho(k, k{+}1, n)$ between the candidate rankings produced by $k$-step and $(k{+}1)$-step scoring at the first decision point on hard-core instances:

\begin{table}[h]
\centering
\begin{tabular}{@{}ccc@{}}
\toprule
$n$ & $\rho(1,2)$ & $\rho(2,3)$ \\
\midrule
7  & 0.249 & 0.960 \\
10 & 0.175 & 0.881 \\
15 & 0.117 & 0.691 \\
18 & 0.105 & 0.392 \\
20 & 0.203 & 0.345 \\
25 & 0.147 & 0.224 \\
\bottomrule
\end{tabular}
\caption{Rank correlation between successive lookahead depths.}
\label{tab:correlation}
\end{table}

Two patterns emerge:

\textbf{1. $k = 1$ to $k = 2$ is a revolution.} $\rho(1,2) \approx 0.1$--$0.2$ at all $n$. The second lookahead level produces a nearly unrelated ranking---it does not refine the first level's assessment but replaces it entirely.

\textbf{2. $k = 2$ to $k = 3$ decays with $n$.} At $n = 7$, $\rho(2,3) = 0.96$---$k = 3$ barely changes anything, and indeed $k = 2$ suffices at this size. At $n = 18$ (where $k = 2$ begins to fail), $\rho(2,3) = 0.39$---significant reshuffling. The correlation drops below ${\sim}0.4$ at precisely the boundary where $k = 2$'s perfect zone ends.

\subsection{Trace Analysis}

At $n = 18$, seed~14 (the first $k = 2$ failure): $k = 2$ scores variable $x_{10} = \mathrm{T}$ at 69 and $x_{18} = \mathrm{F}$ at 68---a near-tie, with $x_{10}$ chosen. This leads to 25 backtracks. Under $k = 3$, $x_{18} = \mathrm{F}$ scores 87, well above $x_{10} = \mathrm{T}$, and the instance is solved in 3 decisions with zero backtracks. The third level of lookahead reveals that $x_{10} = \mathrm{T}$ creates a constrained downstream landscape invisible to 2-step evaluation.

\subsection{Score Ties and Redundancy}

We note a structural property that constrains any proof strategy: the score gap between the best and second-best candidate is typically zero. Large groups of candidates tie for the top score. When $k$ is sufficient, the solver succeeds not because it identifies the unique correct choice, but because \emph{all} tied candidates lead to zero-backtrack solutions. When $k$ is insufficient, ``wrong'' candidates infiltrate the tied group. The fraction of correct candidates in the tied group is ${\sim}80$--$100\%$ when $k$ exceeds the boundary, dropping sharply when $k$ falls below it.

\section{Context and Related Work}\label{sec:context}

\subsection{Algorithmic Thresholds}

Our experiments operate at ratio 4.0, above several well-studied barriers:

\begin{table}[h]
\centering
\begin{tabular}{@{}cl@{}}
\toprule
Ratio & Significance \\
\midrule
${\sim}3.52$ & Best known rigorous polynomial algorithm~\citep{cojaoghlan2010better} \\
${\sim}3.86$ & Clustering/shattering threshold~\citep{achlioptas2008algorithmic} \\
4.0   & \textbf{Our experimental ratio} \\
${\sim}4.15$ & Condensation threshold \\
${\sim}4.267$ & Satisfiability threshold~\citep{dingsly2015proof} \\
\bottomrule
\end{tabular}
\caption{Key thresholds for random 3-SAT.}
\label{tab:thresholds}
\end{table}

Above ratio ${\sim}3.86$, the solution space shatters into exponentially many well-separated clusters~\citep{achlioptas2008algorithmic}, and belief propagation guided decimation provably fails~\citep{cojaoghlan2011belief}. That $k$-step lookahead achieves zero backtracks at ratio 4.0 is notable because it suggests the lookahead mechanism navigates the clustered landscape---but we have no explanation for why.

\subsection{Existing Lookahead Solvers}

Production lookahead solvers (march, kcnfs, OKsolver~\citep{heule2009lookahead}) use related techniques---failed-literal detection, autarky detection---but do not recursively score to arbitrary depth $k$. Our $k$-step model is closer to minimax game-tree search. Whether the structural insights here transfer to practical solver design is an open question.

\subsection{Lower Bounds}

\citet{alekhnovich2011lower} proved exponential lower bounds for DPLL on random 3-SAT. These apply to \emph{refutation} of unsatisfiable instances via tree-like resolution---a different computational task from finding solutions in satisfiable instances. They do not directly constrain our setting, but we note that the distinction between search and refutation is well-understood, and we make no claims about refutation complexity.

Survey Propagation~\citep{mezard2002random} works empirically near the satisfiability threshold but has no polynomial-time guarantee and uses a fundamentally different message-passing approach.

\section{Limitations}\label{sec:limitations}

We want to be explicit about what this paper does and does not establish.

\begin{enumerate}
    \item \textbf{Three data points.} The scaling law rests on three exact measurements. The growth factor is decelerating ($3.13\times$ to ${\geq}2.13\times$). We cannot determine from three points whether it stabilizes, continues to decelerate, or collapses. An exact $k = 5$ data point would significantly strengthen or weaken the case.

    \item \textbf{Small scale.} $n = 180$ is tiny by modern SAT standards ($n > 10^6$). The qualitative character of hard-core instances may change at larger $n$ in ways our measurements cannot anticipate.

    \item \textbf{Weak baselines.} The hard core is defined relative to JW and adaptive polarity---heuristics from the 1990s. A modern CDCL solver (CaDiCaL, Kissat) may solve these instances easily, which would reframe the contribution. We plan this comparison as immediate future work.

    \item \textbf{Beam calibration.} The correction factor (1.237) is measured at one $k$ value. All beam boundaries are lower bounds with unknown tightness.

    \item \textbf{Random instances only.} Industrial/structured SAT instances have different constraint graph structure. Nothing here applies to non-random formulas without further investigation.

    \item \textbf{No proof.} We report an empirical pattern and a structural correlation. We do not prove that $k = O(\log n)$ suffices asymptotically, nor that the diameter connection is causal rather than coincidental.
\end{enumerate}

\section{Open Questions and Future Directions}\label{sec:future}

The diameter--lookahead connection raises several questions we cannot currently answer:

\textbf{Why does diameter-depth lookahead suffice?} The natural conjecture is that when the scorer's information radius exceeds the graph diameter, it ``sees'' all constraints relevant to each decision. But the score-gap analysis (Section~\ref{sec:reshuffle}) shows the mechanism is not simple selection of the unique correct candidate---it is elimination of incorrect candidates via redundancy. Formalizing why redundancy increases with lookahead depth, and whether it connects to the clause-to-variable ratio, would be the most valuable theoretical contribution.

\textbf{Does the growth factor stabilize?} Our three exact points show deceleration: $3.13\times$ then ${\geq}2.13\times$. If it stabilizes above 1, the growth is superlinear regardless of the precise base. If it converges to 1, the growth is eventually linear and the required depth grows as $k = O(n)$. Exact $k = 5$ data---feasible with GPU-accelerated scoring or smarter pruning---would be decisive.

\textbf{Does constant beam width suffice?} Empirically, $B = 3$--$8$ works through $n = 160$. If constant $B$ suffices asymptotically, the per-decision cost with $k = O(\log n)$ is $O(n \cdot B^{O(\log n)}) = O(n^{1 + O(1)})$. If $B$ must grow with $n$, the cost is higher. We have no theoretical basis for predicting either outcome.

\textbf{How does this interact with the Overlap Gap Property?} The OGP framework~\citep{gamarnik2021overlap} constrains ``local'' algorithms on random structures. Whether $k$-step lookahead with $k = O(\log n)$ qualifies as local in the OGP sense---given that it sees the entire constraint graph at the diameter---is unclear and may be an interesting theoretical question independent of our specific results.

\textbf{What do modern solvers achieve on these instances?} A comparison with CaDiCaL or Kissat on the same hard-core instances would establish whether the $k$-step lookahead captures information unavailable to CDCL, or whether CDCL's clause learning achieves comparable backtrack reduction through different means. This is the most important near-term experiment.

\section{Conclusion}

We have measured the zero-backtrack boundary of $k$-step lookahead DPLL on random 3-SAT and found that it grows superlinearly with $k$, tracking the constraint graph diameter. Three exact data points and four beam-estimated lower bounds are consistent with $k = O(\log n)$, though the data is insufficient to establish this asymptotically. The rank-correlation analysis reveals that deeper lookahead operates by reshuffling candidate rankings rather than refining them, with the degree of reshuffling increasing precisely at the point where shallower lookahead fails.

These results hold at ratio 4.0---above the clustering threshold where the solution space shatters and known polynomial algorithms fail---suggesting that multi-level lookahead accesses structural information unavailable to single-pass heuristics. Whether this is a route to efficient algorithms for hard random SAT, or an artifact of small instance sizes, remains to be determined.

The primary contribution of this paper is the empirical observation that required lookahead depth tracks constraint graph diameter. We hope this connection motivates further investigation, both empirical and theoretical, into the information-theoretic role of lookahead depth in constraint satisfaction.

\subsection*{Reproducibility}

All solver source code and instance generators are available at:

\begin{center}
\url{https://github.com/invisiblemonsters/coffinhead-conjecture}
\end{center}

Key reproducible results:
\begin{itemize}
    \item $k=2$ first failure: $n=18$, seed~14, 11~backtracks.
    \item $k=3$ first failure: $n=48$, seed~27, 1506~backtracks.
    \item $k=3$ exact boundary: $n=47$, 30/30 hard core instances perfect.
    \item $k=4$ exact: $n=100$, all hard core instances perfect (parallel, 16 threads).
\end{itemize}
Instance generation uses xorshift64 with the seed as initial state, generating clauses by sampling 3 distinct variables per clause with random polarity.

\subsection*{Acknowledgments}

Computational experiments performed on consumer hardware (WSL2 Ubuntu 24.04).

\bibliographystyle{plainnat}
\bibliography{references}

\end{document}
